% !TEX program = xelatex
\documentclass[11pt]{article}

% ---------- Packages ----------
\usepackage[letterpaper, margin=1in]{geometry}
\usepackage{titlesec}
\usepackage{enumitem}
\usepackage{xcolor}
\usepackage{hyperref}
\usepackage{fontspec}
\usepackage{ragged2e}
\usepackage{setspace}
\usepackage{fontawesome5}
\setstretch{1.1}

% ---------- Colors ----------
\definecolor{accent}{HTML}{1A3A5A}
\definecolor{linkblue}{HTML}{004080}

% ---------- Hyperref ----------
\hypersetup{colorlinks=true, urlcolor=linkblue, linkcolor=linkblue}

% ---------- Section formatting ----------
\titleformat{\section}
  {\color{black}\large\bfseries}
  {}{0em}{\raisebox{0.2ex}{\rule{0.9in}{6pt}}\hspace{0.7em}}
\titlespacing{\section}{0pt}{10pt}{5pt}

% ---------- List spacing ----------
\setlist[itemize]{leftmargin=*, topsep=3pt, itemsep=2pt, parsep=0pt}
\setlength{\parindent}{0pt}
\pagestyle{empty}

% ---------- Entry command ----------
\newcommand{\entry}[2]{%
  \textbf{#1}\hfill{#2}\par\vspace{2pt}%
}

\begin{document}

% ===================== HEADER =====================
\noindent
\begin{minipage}[t]{0.52\textwidth}
  \vspace{0pt}
  {\fontsize{18}{21}\selectfont\bfseries\mbox{Mohammad Amin Shirazian}}\\[5pt]
  {\large\itshape (Amin)}\\[7pt]
  {\Large\itshape Curriculum Vitae}\\[3pt]
  {\footnotesize\href{https://aminshz.github.io/my-website/resume/resume.pdf}{Latest version}}
\end{minipage}%
\hfill
\begin{minipage}[t]{0.40\textwidth}
  \vspace{0pt}
  \raggedleft\small
  \textit{Toronto Metropolitan University}\\
  {\scriptsize 350 Victoria Street, Toronto, ON M5B 2K3, Canada}\\[3pt]
  \faEnvelope\enspace\href{mailto:amin.shirazian.work@gmail.com}{amin.shirazian.work@gmail.com}\\
  \faGlobe\enspace\href{https://aminshz.github.io/my-website}{aminshz.github.io/my-website}\\[3pt]
  {\scriptsize\mbox{Iranian - Canadian Open Work Permit, exp. 2030}}\\[1pt]
  \textit{Updated: September 2026}
\end{minipage}

\vspace{0.8em}

% ===================== RESEARCH FIELDS =====================
\section{Research Fields}

\begin{tabular}{@{}ll@{}}
  Primary: & Macroeconomics, Health Economics \\
  Secondary: & Public Economics, Political Economy
\end{tabular}

% ===================== REFERENCES =====================
\section{References}

\begin{tabular}{p{7cm} p{7cm}}
  \textbf{Keyvan Eslami} & \textbf{Claustre Bajona} \\
  Department of Economics & Department of Economics \\
  Toronto Metropolitan University & Toronto Metropolitan University \\
  \href{mailto:k.eslami@torontomu.ca}{k.eslami@torontomu.ca} & \href{mailto:cbajona@torontomu.ca}{cbajona@torontomu.ca} \\[10pt]

  \textbf{Kevin Fawcett} & \textbf{Haomiao Yu} \\
  Department of Economics & Department of Economics \\
  Toronto Metropolitan University & Toronto Metropolitan University \\
  \href{mailto:fawcettk@torontomu.ca}{fawcettk@torontomu.ca} & \href{mailto:haomiao@torontomu.ca}{haomiao@torontomu.ca} \\
\end{tabular}

% ===================== EDUCATION =====================
\section{Education}

\begin{tabular}{@{}p{1.25in} p{5in}@{}}
  2020--2026 & \textbf{Ph.D. in Economics}, \textit{Toronto Metropolitan University}. \\
  & {\footnotesize Full Scholarship; GPA: A+} \\
  & {\footnotesize Thesis: \textit{``Essays on Health Shocks and Other Precautionary Saving Motives''}} \\
  & {\footnotesize Top Second-Year PhD Student, Toronto Metropolitan University (2021)} \\
  \\[6pt]

  2018--2020 & \textbf{M.A. in Economics}, \textit{Tehran Institute for Advanced Studies}. \\
  & {\footnotesize Full Scholarship; GPA: A} \\
  & {\footnotesize Ranked 68th nationally in Iran's MA Economics entrance examination} \\[6pt]

  2011--2017 & \textbf{B.A. in Economics}, \textit{University of Tehran}. \\
  & {\footnotesize Highly competitive, tuition-free admission through Iran's national entrance examination; GPA: B} \\
\end{tabular}

% ===================== SKILLS =====================
\vspace{-4pt}
\section{Skills}

{\footnotesize\noindent\textbf{Statistical / data:} Stata, R, Python (Selenium, BeautifulSoup) \,$\cdot$\,
\textbf{Computational:} MATLAB, C++ \,$\cdot$\,
\textbf{Other:} \LaTeX{}, Git, HTML/CSS \,$\cdot$\,
\textbf{Languages:} English (fluent), Persian (native)}

\pagebreak
\small

% ===================== JOB MARKET PAPER =====================
\section{Job Market Paper}

\textbf{Marginal Utility Shocks and the Precautionary Saving Puzzle}
\hfill
{\footnotesize\href{https://aminshz.github.io/my-website/paper1.pdf}{latest version}}

{\footnotesize\textit{Documenting that both saving and medical spending accelerate at the top, the paper derives a sufficient condition for saving to remain positive at arbitrarily high wealth and shows that the resulting stationary wealth distribution has Pareto tails, the first such result from precautionary saving alone. Calibrated to the data, the mechanism accounts for roughly 11 percent of observed saving in the top wealth decile.}}

% ===================== WORK IN PROGRESS =====================
\section{Work in Progress}

\entry{Means-Tested Insurance and Pareto-Tailed Wealth Distribution}{In Progress}
\textit{This paper studies how means-tested healthcare programs affect household wealth accumulation.
The usual view emphasizes that households may decumulate wealth to qualify for coverage. We identify
a countervailing dynamic force: in each period, households choose whether to reduce their assets and
become eligible or to remain ineligible and preserve their wealth. Repeated decisions to remain
ineligible can generate sustained wealth accumulation and contribute to wealth concentration at the
top of the distribution. Evaluating this eligibility choice empirically requires total medical spending,
including the insurance component. Because this measure is unavailable in the Health and Retirement
Study, the paper imputes it using the Medical Expenditure Panel Survey (MEPS). It develops conditions
for consistency of the imputation and of OLS estimates using the imputed dependent variable. A stylized model shows how
an asset threshold for means-tested insurance can generate strategic asset decumulation, linking
insurance eligibility to household saving behavior.}

\vspace{1em}

\entry{Online Innovation, Market Entry and Competition in Remote Indigenous Communities}{In Progress}
\textit{With Nicholas Li and Yasir Mudathir.}

\vspace{1em}

\entry{Health Production Function and Longevity}{In Progress}
\textit{This project studies the health production function by identifying which health products and
medical technologies improve survival, how effectively private treatment supplements available care,
and how rapidly returns to additional medical spending diminish. It examines the wealth threshold at
which households begin to demand these products and the characteristics beyond wealth that shape the
demand for longevity, while addressing the two-way relationship between longevity preferences and
wealth accumulation.}

\vspace{1em}

\entry{Optimal Insurance Against Mortality Shocks}{In Progress}
\textit{This project extends the current framework to a general-equilibrium model with bequest motives
and family dynamics. The model captures how households share health risks, provide informal care,
transfer wealth across generations, and save for medical needs and bequests. It evaluates how much of
saving across the wealth distribution these motives explain, whether household choices under the
existing healthcare and insurance system are efficient, and when incomplete insurance markets create
a rationale for government intervention.}

\vspace{1em}

\entry{Gender and Longevity: A Retirement Saving Puzzle}{MA Thesis}
\textit{This paper presents an explanation to the excess sensitivity of consumption to the drop in income
around retirement, known as the ``retirement saving puzzle.'' We confirm this empirical puzzle also
prevails among Iranian families. That is, an unpredicted drop occurs in households' non-durable
expenditures around the age of 65, which cannot be resolved even by decomposing household heads
according to their labor status. Then, by taking into account the diversity in the gender composition
of cohorts, we demonstrate that the gap between predicted and actual non-durable expenditures is no
longer apparent.}

\newpage
\titlespacing{\section}{0pt}{7pt}{4pt}

% ===================== TEACHING =====================
\section{Teaching \& Work Experience}

\textbf{Teaching}\\[5pt]

\begin{tabular}{p{2.5cm} p{15cm}}
  Summer 2024 -- Present &
    \textbf{Contract Lecturer}, Department of Economics, Toronto Metropolitan University.
    \\[-1pt]
  & {\scriptsize Fall 2026: CECN230 -- \textit{Mathematics for Economics} (enrollment TBD)} \\
  & {\scriptsize Summer 2026: CECN600 -- \textit{Intermediate Macroeconomics II} (20 students)} \\
  & {\scriptsize Spring 2026: CECN603 -- \textit{Economic Issues in Globalization} (53 students)} \\
  & {\scriptsize Winter 2026: CECN301 -- \textit{Intermediate Macroeconomics I} (87 students)} \\
  & {\scriptsize Fall 2025: CECN301 -- \textit{Intermediate Macroeconomics I} (38 students)} \\
  & {\scriptsize Summer 2025: CECN603 -- \textit{Economic Issues in Globalization} (44 students)} \\
  & {\scriptsize Winter 2025: CECN301 -- \textit{Intermediate Macroeconomics I} (61 students)} \\
  & {\scriptsize Summer 2024: CECN204 -- \textit{Introductory Macroeconomics} (59 students)} \\
\end{tabular}

\vspace{9pt}
\textbf{Graduate Assistant}\\[3pt]

\begin{tabular}{p{2.5cm} p{15cm}}
  2018 -- Fall 2026 & \textbf{Teaching Assistant} \\[2pt]
  & Department of Economics, Toronto Metropolitan University. \\
  &
    Taught \textit{Adv.\ Macroeconomics I}, \textit{Mathematical Economics},
    \textit{Advanced Microeconomics II}; led PhD tutorials \\[6pt]

  &
    Tehran Institute for Advanced Studies / Khatam University.
    Taught \textit{Advanced Macroeconomics}, \textit{Advanced Financial Economics}; led MA tutorials \\[3pt]
  & {\scriptsize Graduate Assistant Recognition Award, Toronto Metropolitan University (2022--2023)} \\[9pt]

  2022--2024 & \textbf{Research Assistant} \\[2pt]
  & Department of Economics, Toronto Metropolitan University \\
  & {\scriptsize Prof. Keyvan Eslami -- \textit{Imputing Total Medical Spending}: Combined HRS and MEPS data to impute total medical spending for econometric analysis.} \\[3pt]
  & Department of Economics, University of Toronto \\
  & {\scriptsize Prof. Nicholas Li -- \textit{Nutrition North Price Data}: Built a web-scraping workflow to collect and organize prices from Nutrition North program providers' websites.} \\[3pt]
  & Tehran Institute for Advanced Studies, Tehran, Iran \\
  & {\scriptsize Prof. Farshad Hagh-Panah -- \textit{Household Income and Expenditure Survey Database}: Collected, streamlined, and cleaned HIES data and constructed a database of households classified by the gender of the household head.} \\
\end{tabular}

% ===================== PROFESSIONAL EXPERIENCE =====================
\vspace{9pt}
\textbf{Union Experience}\\[3pt]

\begin{tabular}{p{2.5cm} p{15cm}}
  2021--2023 & \textbf{Economics Graduate Students Association}, Toronto Metropolitan University \\
  & \textit{President} (2021--2022); \textit{VP Communication} (2022--2023) \\
  & {\scriptsize Student-association leadership focused on planning events, coordinating funding with the student union and faculty, and managing teams, invitations, and logistics.} \\[4pt]
  2021--2022 & \textbf{Union Governance -- Faculty of Arts Director}, Toronto Metropolitan Graduate Students' Union (TMGSU) \\
  & {\scriptsize Represented the Faculty of Arts on the union board; reviewed bylaws, attended governance meetings, helped develop the newly established union's structure, and contributed to funding decisions for health insurance and graduate-student travel.} \\
\end{tabular}

\normalsize

\vfill

\end{document}
